\subsubsection{伴随边值问题}\label{subsec:adjoint-boundary-value-problem}

本小节始终设 $L(\bm{x},D)$ 为 $2m$ 阶标量椭圆微分算子, $B_j(\bm{x},D)$, $j=1,\ldots,m$, 为满足补条件的边界算子. 伴随问题的一般理论要求 $\Omega$ 和系数满足相当严格的正则性假设; 为简单起见, 假设它们属于 $C^\infty$ 类且 $\Omega$ 有界. 以下始终作这些假设. 还将附加假设 $B_j$ 正规. 这一性质定义如下.

\begin{Definition}
\label{def:normal-boundary-operators}
若边界算子 $B_j(\bm{x},D)$ 的阶数 $m_j$ 两两不同且不超过 $2m-1$, 并且 $B_j$ 的最高阶项含有纯法向导数, 即对每个 $\bm{x}\in\partial\Omega$ 都有 $B_j^p(\bm{x},\bm{n})\ne0$ (这里 $\bm{n}$ 是 $\partial\Omega$ 的单位法向量), 则称这些边界算子正规.
\end{Definition}

$B_j$ 的阶数只覆盖从 $0$ 到 $2m-1$ 的一半取值. 可以添入边界算子 $S_j$, $j=1,\ldots,m$, 补齐缺少的阶数. 显然可以使扩充后的边界算子组仍满足正规性条件; 只需取 $S_j$ 为 $\partial/\partial n$ 的适当幂即可. 作如下定义.

\begin{Definition}
\label{def:dirichlet-boundary-system}
若边界算子 $F_j(\bm{x},D)$, $j=1,\ldots,p$, 的阶数 $m_j$ 覆盖从零到 $p-1$ 的全部取值, 并且 $F_j$ 的最高阶项含有纯法向导数, 即对每个 $\bm{x}\in\partial\Omega$ 都有 $F_j^p(\bm{x},\bm{n})\ne0$ (这里 $\bm{n}$ 是 $\partial\Omega$ 的单位法向量), 则称这些边界算子为 $p$ 阶 Dirichlet 系统.
\end{Definition}

有以下引理.

\begin{Lemma}
\label{lem:dirichlet-system-tangential-triangular-reduction}
设 $F_i(\bm{x},D)$, $i=1,\ldots,p$, 为 Dirichlet 系统, 并设 $F_i$ 的阶数为 $i-1$. 则存在 $i-j$ 阶切向微分算子 $\Phi_{ij}(\bm{x},D)$ 和 $\Psi_{ij}(\bm{x},D)$, 使得
\begin{equation}\label{eq:dirichlet-system-tangential-triangular-reduction}
\begin{aligned}
F_i(\bm{x},D)&=\sum_{j=0}^i\Phi_{ij}(\bm{x},D)\frac{\partial^{j-1}}{\partial n^{j-1}},\\
\frac{\partial^{i-1}}{\partial n^{i-1}}&=\sum_{j=0}^i\Psi_{ij}(\bm{x},D)F_j(\bm{x},D).
\end{aligned}
\end{equation}

\end{Lemma}

$\Phi_{ij}$ 的存在性由定义显然可知. 然后将 $\Phi_{ij}$ 构成的三角矩阵求逆, 即可得到 $\Psi_{ij}$. 证明的细节留作练习; 见 Problem~\ref{ex:dirichlet-system-triangular-reduction-details}.

\begin{Corollary}
\label{cor:dirichlet-system-prescribed-traces}
设 $F_i$, $i=1,\ldots,2m$, 为 Dirichlet 系统, 用 $m_i$ 表示 $F_i$ 的阶数. 给定 $g_i\in H^{2m+k-m_i-1/2}(\partial\Omega)$. 则存在 $u\in H^{2m+k}(\Omega)$, 使得在 $\partial\Omega$ 上有 $F_iu=g_i$.
\end{Corollary}

证明由前一引理和 Theorem~\ref{thm:domain-higher-normal-trace} 立即得到.
